\chapter{Einleitung}
\label{chap:einleitung}

\section{Motivation}
\label{sec:motivation}

Bühnenbeleuchtung ist ein zentraler Bestandteil moderner Theater-, Konzert- und
Veranstaltungsproduktionen. Durch gezielte Lichtgestaltung können Atmosphäre,
Stimmungen und visuelle Effekte erzeugt werden, die wesentlich zur Wahrnehmung
einer Aufführung durch das Publikum beitragen. Neben der rein funktionalen
Ausleuchtung der Bühne ermöglicht die Lichttechnik somit eine kreative und
dramaturgische Gestaltung des gesamten Bühnengeschehens.

Zur Steuerung professioneller Beleuchtungssysteme hat sich in der
Veranstaltungstechnik das Protokoll DMX (Digital Multiplex) als De-facto-Standard
etabliert. DMX erlaubt die zentrale Steuerung einer Vielzahl von
Beleuchtungsgeräten über ein standardisiertes Kommunikationsprotokoll und bietet
eine zuverlässige sowie flexible Möglichkeit zur Echtzeitsteuerung komplexer
Lichtinstallationen.

Die Signalführung dieses Standards ist jedoch kabelgebunden: Jedes
Beleuchtungsgerät muss über eine eigene Leitung in eine durchgehende DMX-Kette
eingebunden werden. Vor allem bei temporären Installationen, mobilen Bühnen oder
häufig wechselnden Aufbauten wird die dafür erforderliche Verkabelung zu einem
wesentlichen Komplexitätsfaktor. Vor diesem Hintergrund entsteht ein Interesse an
flexibleren Lösungen, die den Verkabelungsaufwand bei gleichbleibender
Funktionalität verringern.

\section{Problemstellung}
\label{sec:problemstellung}

Die kabelgebundene Signalübertragung bietet eine hohe Zuverlässigkeit, führt in
der Praxis jedoch zu drei Nachteilen, die den Einsatz insbesondere außerhalb
fester Installationen erschweren.

Der erste betrifft Zeit und Personal. Die Signalverkabelung erfordert eine
sorgfältige Planung sowie die manuelle Verbindung sämtlicher Beleuchtungsgeräte
innerhalb der DMX-Kette. Bei größeren Installationen entsteht dadurch ein
erheblicher Zeitaufwand beim Auf- und Abbau, der zudem Personal mit
entsprechender Erfahrung im Umgang mit Veranstaltungstechnik voraussetzt.

Der zweite Nachteil ist materieller Art. Kabel und Verbindungselemente
verursachen Kosten für Anschaffung, Wartung und Transport und müssen in
ausreichender Länge und Anzahl vorgehalten werden. Mit der Zahl der Verbindungen
steigt darüber hinaus die Zahl möglicher Fehlerquellen, was die Fehlersuche im
laufenden Betrieb erschwert.

Der dritte Nachteil liegt im verfügbaren Angebot. Zwar existieren bereits
kommerzielle Lösungen zur drahtlosen Übertragung von DMX-Daten, diese sind jedoch
häufig mit hohen Anschaffungskosten verbunden und richten sich primär an
professionelle Großproduktionen. Für kleinere Produktionen, experimentelle
Bühneninstallationen oder flexible Lichtsysteme stehen daher nur begrenzt
kostengünstige und zugleich leistungsfähige Lösungen zur Verfügung.

\section{Zielsetzung}
\label{sec:zielsetzung}

Ziel dieser Arbeit ist die Entwicklung eines Beleuchtungssystems für
Veranstaltungen, das DMX-Steuersignale drahtlos an mehrere Leuchteinheiten
überträgt. Diese werden in Form von LED-Leuchtröhren realisiert, die die
empfangenen Steuerdaten unmittelbar in Lichtsignale umsetzen. Ein zentrales
Element bildet eine sogenannte Bridge, welche die Steuerdaten aus einer
bestehenden Lichtinfrastruktur entgegennimmt und an die Leuchtröhren verteilt.
Die Arbeit umfasst dabei sowohl die Hardware als auch die Software beider
Gerätetypen.

Ein wesentliches Entwicklungsziel ist die Gewährleistung einer stabilen
Datenübertragung mit hinreichender Aktualisierungsrate und zeitlicher Konstanz,
sodass die zeitkritischen Anforderungen der Lichtsteuerung erfüllt werden. Dies
schließt die robuste Funktion unter praxisnahen Bedingungen mit potenziellen
Funkstörungen ausdrücklich ein.

Darüber hinaus soll der Nachbau des Systems niedrigschwellig möglich sein.
Vorgesehen sind deshalb ausschließlich frei zugängliche Komponenten, sodass auch
Anwenderinnen und Anwender mit geringen Vorkenntnissen den Aufbau bewältigen
können. Diesem Ziel folgt auch die Ausrichtung als Open-Source-Projekt: Neben der
Software werden die Platinenlayouts der Elektronikbaugruppen sowie die 3D-Modelle
der Gehäuse veröffentlicht. Damit sollen Transparenz geschaffen, die
Nachvollziehbarkeit technischer Entscheidungen verbessert und die
Weiterentwicklung durch Dritte erleichtert werden.

\section{Wissenschaftliche und praktische Relevanz}
\label{sec:wissenschaftliche-und-praktische-relevanz}

Die Kommunikation in verteilten eingebetteten Systemen ist ein relevantes
Forschungsfeld innerhalb der Informatik. Zuverlässigkeit, Latenzverhalten,
Skalierbarkeit und Robustheit unter realen Einsatzbedingungen stellen dabei
zentrale Fragestellungen dar. Das entwickelte System realisiert ein solches
verteiltes eingebettetes Echtzeitsystem zur Steuerung von Aktoren und erlaubt es
damit, Entwurfsentscheidungen, Leistungsgrenzen und Zielkonflikte zwischen
Flexibilität, Komplexität und Übertragungsqualität systematisch zu untersuchen.
Von besonderem Interesse ist die Frage, welche Zuverlässigkeit eine
unidirektionale Funkstrecke ohne Zustellbestätigung unter realen
Umgebungsbedingungen erreicht.

Praktisch adressiert die Arbeit genau jenen Anwendungsbereich, für den nach
Abschnitt~\ref{sec:problemstellung} bislang keine geeignete Lösung verfügbar ist.
Da sämtliche Entwurfsdaten offengelegt werden, ist das Ergebnis nicht nur
nachvollziehbar, sondern auch unmittelbar nachnutzbar.

\section{Aufbau der Arbeit}
\label{sec:aufbau-der-arbeit}

Kapitel~\ref{chap:grundlagen} stellt die theoretischen Grundlagen zu
eingebetteten Systemen, den verwendeten Kommunikationsprotokollen und der
drahtlosen Übertragung vor. Kapitel~\ref{chap:stand-der-technik} ordnet die
Arbeit in den Stand der Technik ein, indem kommerzielle Systeme und verwandte
wissenschaftliche Arbeiten betrachtet und die bestehenden Lücken benannt werden.

Auf dieser Grundlage leitet Kapitel~\ref{chap:anforderungsanalyse} die
funktionalen und nicht-funktionalen Anforderungen ab.
Kapitel~\ref{chap:systemarchitektur} beschreibt daraufhin die Architektur des
Systems einschließlich der Hardware- und Softwarekomponenten von Bridge und
Leuchtröhren, während Kapitel~\ref{chap:implementierung} deren konkrete Umsetzung
darlegt.

Kapitel~\ref{chap:evaluierung} untersucht das entwickelte System experimentell.
Erhoben werden Messreihen zur Verbindungsstabilität, zur Aktualisierungsrate und
zur Zuverlässigkeit der Übertragung, die anschließend den Anforderungen
gegenübergestellt werden. Kapitel~\ref{chap:diskussion-und-ausblick} bewertet die
Ergebnisse, zieht ein Fazit und zeigt Ansätze für zukünftige Weiterentwicklungen
auf.
